\section{Highly supersymmetric Lorentzian spin manifolds}
\label{sec:high-susy-spin-mfld}

In this section, we invert the perspective of Section~\ref{sec:killing-spinors-superalgebras}: rather than constructing a filtered subdeformation of a flat model $\fs$ from sections on a spin manifold (namely the Killing superalgebra), we seek to begin with such a filtered subdeformation and construct a spin manifold with admissible connection on which the elements of the algebra can be realised as sections, in doing so providing a partial converse to Theorem~\ref{thm:killing-algebra-filtered}. We restrict ourselves to the highly supersymmetric Lorentzian case, in which the homogeneity theorem (see Theorem~\ref{thm:homogeneity}) provides significant simplifications as it did in Section~\ref{sec:filtered-def-poincare} (although see Remark~\ref{rem:relaxing-homogeneity-assumption}), and our results will ultimately justify the approach of that section by demonstrating that our notion of geometric realisability (see Definition~\ref{def:geom-real}) is a good one. We will essentially generalise the results of \cite[Section~3.3]{Figueroa-OFarrill2017_1}. We make extensive use of the conventions and results related to homogeneous spaces discussed in Appendix~\ref{sec:homog-spaces}. We begin by establishing how we will deal with spinors on homogeneous manifolds.

\subsection{Homogeneous spin structures}
\label{sec:homogeneous-spin-mfld}

Spin structures on homogeneous manifolds were first treated systematically by B\"ar \cite{Bar1992} and by Cahen et al.\ \cite{Cahen1993}; our present treatment follows \cite{Figueroa-OFarrill2017_1}.

Let $(G,H,\eta)$ be a metric Klein pair (of any signature), let $(M=G/H,g)$ be the associated homogeneous space, which we assume to be oriented, and let $V=T_oM$ where $o=H$. Suppose that there exists a lift of the isotropy representation $H\to\SO(V)$ to $\Spin(V)$; that is, a Lie group morphism $H\to\Spin(V)$ making the following diagram commute
\[
\begin{tikzcd}
	&	& \Spin(V) \ar[d]	\\
	& H	\ar[r] \ar[ur,dashed]& \SO(V).
\end{tikzcd}
\]
Then the canonical map $\Spin(V)\to\SO(V)$ induces a $\Spin(V)$-equivariant bundle map
\[
	G\times_H\Spin(V)\to G\times_H\SO(V) \cong F_{\rm SO},
\]
which is a spin structure on $M$, where we recall that $F_{\rm SO}$ is the special orthonormal frame bundle. In order to define the isomorphism we have implicitly chosen an orthonormal frame over $o$ (as discussed in Appendix~\ref{sec:homog-spaces-principal-bundles}). We call this the \emph{homogeneous spin structure associated to the lift $H\to\Spin(V)$}. Note that $G$ (left) acts naturally on both of these bundles compatibly with the right action by the spin group, and that the bundle isomorphism is $G$-equivariant.

Using the isotropy representation and related representations of $H$, many natural bundles on $M$, particularly those associated to $F_{\rm SO}$ via representations of $\SO(V)$, can be viewed as associated bundles to the $H$-principal bundle $G\to M$. Homogeneous spin structures allow us to do the same with representations of the spin group; for example, if $S$ is a spinor module of~$\Spin(V)$, there is a $G$-equivariant isomorphism of vector bundles
\[
	\Sbundle := (G\times_H\Spin(V))\times_{\Spin(V)} S \cong G\times_H S.
\]

If $G$ is simply connected, the construction above classifies the spin structures on $(M,g)$ in definite signature.

\begin{Lemma}[\cite{Bar1992,Cahen1993}]\label{lemma:homogeneous-spin}
	If $(M,g)$ is a homogeneous Riemannian $G$-space where $G$ is simply connected, the equivalence classes of spin structures on $M$ are in one-to-one correspondence with spin lifts of the isotropy representation.
\end{Lemma}

The proof of this result relies on the ability to lift the action of $G$ on $F_{\rm SO}$ induced by the action on $M$ to an action on an arbitrary spin structure $P\to M$ which can then be localised at $o$ to find a lift of the isotropy representation. In definite signature, this follows from the fact that the spin group is connected and simply connected (for $\dim V \geq 3$), but it fails in indefinite signature due to the non-trivial topology of the spin group; in signature $(p\geq 1,q\geq 1)$, the spin group is disconnected, and its connected component is not simply connected for $(p\geq 2,q\geq 2)$. However, in the case that $M$ is a \emph{reductive} $G$-space, the above still holds in arbitrary signature~\cite{Alekseevsky2019,Cahen1993}.

To avoid these issues, one can simply hypothesise that the action of $G$ lifts to the spin structure; we then say that we have a \emph{homogeneous spin structure}. Such structures are in one-to-one correspondence with lifts of the isotropy representation.

\subsection{Reconstruction of highly supersymmetric backgrounds}
\label{sec:high-susy-spin-mfld-homogeneous}

Let us now turn directly to the question of constructing a background geometry on which a~filtered deformation of a Poincar\'e superalgebra can be realised. Let $(V,\eta)$ be a Lorentzian inner product space, $S$ a (possibly $N$-extended) spinor module of $\Spin(V)$, $\kappa\colon \Odot^2 S\to V$ a symmetric and causal Dirac current and $\fs$ the Poincar\'e superalgebra defined by this data (see Definition~\ref{def:flat-model-alg}). Our ultimate goal is to generalise \cite[Theorem~13]{Figueroa-OFarrill2017_1}, which says that geometrically realisable highly supersymmetric filtered subdeformations of the 11-dimensional Poincar\'e superalgebra are subalgebras of Killing superalgebras. We will find some potential obstructions to doing this in full generality, but the obstructions themselves will prove instructive.

\subsubsection{Super Harish-Chandra pairs}

We begin by introducing what can be viewed as a Lie supergroup associated to a Lie superalgebra. In a superspace formalism, the supersymmetric backgrounds we will construct could be considered as homogeneous spaces for such a supergroup. We will not fully use such a formalism, but the following definition will be useful for us. See \cite{Santi2010} and the references therein for more on the superspace perspective.

\begin{Definition}[super Harish-Chandra pair]\label{def:super-Harish-Chandra}
	A \emph{super Harish-Chandra pair} is a pair $(G_{\overline0},\fg)$ where $\fg$ is a finite-dimensional Lie superalgebra and $G_{\overline0}$ is a connected Lie group integrating $\fg_{\overline0}$ equipped with a representation $\Ad\colon G_{\overline0}\to\GL(\fg)$ (which we call the \emph{adjoint representation} of the pair) integrating the restriction $\ad\colon\fg_{\overline0}\to\fgl(\fg)$ of the adjoint representation of $\fg$ to $\fg_{\overline0}$.
\end{Definition}

Note that if $(G_{\overline0},\fg)$ is a super Harish-Chandra pair, the adjoint representation $\Ad\colon G_{\overline0}\to\GL(\fg)$ preserves the grading on $\fg$ since $\ad$ does. The induced subrepresentation on $\fg_{\overline0}$ is simply the adjoint representation of $G_{\overline0}$.

Let us now consider two motivating examples.

\begin{Example}
The spin cover $\Spin_0(V)\ltimes V$ of the connected Poincar\'e group $\SO_0(V)\ltimes V$ of course has the Poincar\'e algebra $\fs_{\overline 0}=\fiso(V)$ as its Lie algebra, and equipping it with the action $(a,v)\cdot s = a\cdot s$ on $\fs_{\overline 1}=S$ makes $(\Spin_0(V)\ltimes V,\fs)$ into a super Harish-Chandra pair.
\end{Example}

\begin{Example}
Let $(M,g,D)$ be a connected Lorentzian spin manifold with admissible connection and let $\fK_D$ be its Killing superalgebra (see Definition~\ref{def:killing-spinor}). Let us further assume that this is \emph{highly supersymmetric}, in the sense that $\dim\fS_D>\frac{1}{2}\dim S$ (i.e., $\fK_D$ is highly supersymmetric as a filtered subdeformation of $\fs$). Then, by Theorem~\ref{thm:homogeneity}, the values of the Killing vector fields arising as squares of Killing spinors span the tangent space at every point of $M$. Thus $M$ is \emph{locally homogeneous}; informally, it locally looks like a homogeneous space for a group with Lie algebra $\fV_D$. If we further assume that $M$ is geodesically complete, $\fV_D$ integrates to a connected subgroup $G$ of $\ISO(M,g)$ which acts transitively on $M$ \cite[Theorem~2.2]{Stehney1980}, whence $M$ is in fact homogeneous. By pulling back the action of $G$ to its universal covering group $\widetilde{G}$, it follows that $(M,g)$ is a homogeneous Lorentzian spin manifold for $\widetilde{G}$. Since $\widetilde{G}$ is simply connected, the restriction $\ad\colon\fV_D\to\fgl(\fK_D)$ of the adjoint representation of $\fK_D$ to $\fV_D$ integrates to a~representation $\Ad\colon\widetilde{G}\to\fK_D$. Thus $(\widetilde{G},\fK_D)$ is a super Harish-Chandra pair.
\end{Example}

The second example in particular shows us a way forward: where a Killing superalgebra is highly supersymmetric, the background which supports it is locally homogeneous, and if it is actually homogeneous, one can associate to it a super Harish-Chandra pair containing the Killing superalgebra.

Consider now a highly supersymmetric filtered subdeformation $\fg$ of the Poincar\'e superalgebra~$\fs$, say $\Gr\fg\cong\fa\subseteq\fs$ as graded Lie superalgebras so that $V=\fa_{-2}$, $S':=\fa_{-1}$ and $\fh:=\fa_0$ as usual. In particular, the filtration on $\fg$ is of the form $\fg=\fg^{-2}\supseteq\fg^{-1}\supseteq\fg^0\supseteq\fg^1=0$ and as Lie algebras, $\fh=\Gr_0\fg=\fg^0/\fg^1\cong\fg^0$, so there is a canonical embedding of $\fh$ as the zeroth part of the filtration of $\fg$.\footnote{This can also be seen by considering an explicit presentation of $\fg$ as a filtered deformation of $\fa$.}
Note also that the adjoint action of $\fa^0$ preserves each $\fg^k$ so that $\fg$ is also filtered as an $\fg^0$-module, whence $\Gr\fg$ is also naturally a graded $\fg^0$-module, and this action is the same one induced by the isomorphism $\fh\cong\fg^0$. We will thus identify $\fg^0$ and $\fh$, in particular treating $\fh$ as a subalgebra of $\fg$.

\begin{Lemma}\label{lemma:spin-lift-Ad}
	Let $(G_{\overline0},\fg)$ be a super Harish-Chandra pair where $\fg$ is a highly supersymmetric filtered subdeformation of $\fs$, and suppose that the connected subgroup $H$ of $G_{\overline0}$ generated by~${\fh=\fg^0}$ is closed. Then the inclusion $\fh\hookrightarrow\fso(V)$ integrates to a morphism $H\to\Spin(V)$. In particular, this lifts the isotropy representation and thus defines a homogeneous spin structure for $(G_{\overline0},H,\eta)$.
\end{Lemma}

\begin{proof}
	First note that $V\cong \fg^{-2}/\fg^{-1} \cong \fg_{\overline0}/\fh$ and $S'\cong\fg^{-1}/\fg^0\cong\fg_{\overline1}$ as $\fh$-modules. In particular, since $\eta$ is $\fh$-invariant, $(\fg_{\overline0},\fh,\eta)$ is a metric Lie pair, so $(G_{\overline0},H,\eta)$ is a metric Klein pair, and the inclusion $\fh\hookrightarrow\fso(V)$ is the isotropy representation of $\fh$. Similarly, at the level of groups, the adjoint representation induces a Lie group morphism $\varphi\colon H\to\SO(V)$ integrating $\fh\hookrightarrow\fso(V)$. This lifts to a morphism $\tilde\varphi\colon\Htilde \to \Spin(V)$ where $\pi_H\colon\Htilde\to H$ is the universal cover of $H$. We will show that there exists a diagonal arrow making the diagram
		\[
		\begin{tikzcd}
			& \Htilde \ar[d,"\pi_H"]\ar[r,"\tilde\varphi"]	& \Spin(V) \ar[d,"\pi"]	\\
			& H \ar[r,"\varphi"]\ar[ur,dotted]	& \SO(V)
		\end{tikzcd}
		\]
	commute, which proves the result. Recalling that there is a canonical short exact sequence of Lie groups
$ 1 \longrightarrow \pi_1(H) \longrightarrow\Htilde \longrightarrow H \longrightarrow1$,
	we will show that $\pi_1(H)\subseteq\ker\tilde\varphi$, hence $\tilde\varphi$ factors through $H$, giving the required arrow in the diagram.
	
	Note that the actions of $H$ on $V$ and $S'$ pull back to actions of $\Htilde$ (integrating the actions of~$\fh$) in which $\pi_1(H)$ acts trivially. On the other hand, $\Htilde$ also acts on $V$ and $S$ via the action of~$\Spin(V)$ on those spaces. The two actions of $\Htilde$ on $V$ agree by the commutation of the square diagram above. Since $\tilde\varphi$ integrates $\fh\hookrightarrow \fso(V)$ and the action of $\fh$ (via $\fso(V)$) on $S$ preserves $S'$, so too does the action of $\Htilde$ via $\Spin(V)$. The two actions of $\Htilde$ on $S'$ both integrate the same action of $\fh$ on $S'$, so they agree.
		
	The trivial action of $\pi_1(H)$ on $V$ gives us
$\pi_1(H)\subseteq\ker(\varphi\circ\pi_H)=\ker(\pi\circ\tilde\varphi)
$
	thus $\tilde\varphi(\pi_1(H))\subseteq\ker\pi=\{\pm\1\}$. On the other hand, the trivial action of $\pi_1(H)$ on $S'$ shows that $-\1\notin\tilde\varphi(\pi_1(H))$, thus $\pi_1(H)\subseteq\ker\tilde\varphi$. This shows that $\tilde\varphi$ factors through $H$ as claimed.
\end{proof}

\subsubsection{Reconstruction theorem}

We can now state and prove our final result which is a partial converse to Theorem~\ref{thm:killing-algebra-filtered} and generalisation of \cite[Theorem~13]{Figueroa-OFarrill2017_1}, with some caveats. Since this result essentially allows one to (locally) reconstruct a highly supersymmetric background from its Killing superalgebra, we follow \cite{Figueroa-OFarrill2017_1,Santi2022} and refer to it as the \emph{Reconstruction Theorem}.

\begin{Theorem}[reconstruction of highly supersymmetric background]\label{thm:homog-spin-mfld}	
	Let $\fg$ be a geometrically realisable highly supersymmetric filtered subdeformation of the Poincar\'e superalgebra $\fs$.
	Let $G_{\overline0}$ be the connected and simply connected Lie group corresponding to the Lie algebra $\fg_{\overline0}$ and suppose that the connected subgroup $H$ corresponding to $\fh$ is closed.
	Then there exists a homogeneous spin structure on the Lorentzian homogeneous space $M=G_{\overline0}/H$ with a Dirac current $\kappa$ and a~connection $D$ on the spinor bundle $\Sbundle=G_{\overline 0}\times_H S$ as well as an injective $\ZZ_2$-graded linear map~${\Psi\colon\fg\hookrightarrow\fV_D\oplus\fS_D}$ which restricts to a Lie algebra embedding $\fg_{\overline 0}\hookrightarrow\fV_D$ and which satisfies
	\begin{align}\label{eq:reconstruct-brackets}
		\Psi(\comm{X}{s}) = \eL_{\Psi(X)}\Psi(s),
		\qquad \Psi\qty(\comm{s}{s'}) = \kappa\qty(\Psi(s),\Psi(s'))
	\end{align}
	for all $X\in\fg_{\overline 0}$ and $s,s'\in \fg_{\overline 1}$.
	
	Furthermore, if $\Psi(\fg_{\overline 1})=\fS_D$ or if $D$ is admissible in the sense of Definition~{\rm\ref{def:killing-spinor}} otherwise, we find in particular that $\fg$ embeds into the Killing superalgebra $\fK_D$ of $(M,g,D)$.
\end{Theorem}

\begin{Remark}
	Let us make some comments about this result before giving the proof. First, it is necessary to hypothesise that $H$ is closed; it may be possible to adjust the choice of $\fh$ such that this is the case but we have not been able to show this. We note that this is also implicitly assumed in \cite[Theorem~13]{Figueroa-OFarrill2017_1}; although it is not in the statement of the theorem, it is stated in the preamble to the theorem.
	
	Moreover, the present theorem is very close to saying that \emph{all} geometrically realisable highly supersymmetric subdeformations of $\fs$ embed into Killing superalgebras, but there is an obstruction in that it does not seem to be guaranteed that the connection $D$ on $\Sbundle$ which we will construct in the proof is admissible. The algebraic conditions~(1) and~(2) of Definition~\ref{def:killing-spinor} are satisfied because we work with admissible cocycles which are expressed in terms of normalised cocycles of the whole Poincar\'e superalgebra $\fs$; the difficulty is in showing condition~(3), namely that $\eL_{\kappa_\epsilon}\beta=0$ for all $\epsilon\in\fS_D$; we can only show this for $\epsilon\in\Psi(\fg_{\overline 1})$. On the other hand, by a~dimension count we find that \emph{maximally} supersymmetric subdeformations do in fact embed in Killing superalgebras, since in this case we must have $\Psi(\fg_{\overline 1})=\fS_D$.
	
	To put a positive spin on this issue, it provides us with new motivation to study Killing superalgebras constructed from \emph{constrained} Killing spinors, which we briefly discussed in Section~\ref{sec:alg-kse}; while $D$ might not be guaranteed to satisfy the conditions for admissibility on sections of $\Sbundle$, it \emph{does} satisfy these conditions on sections of the sub-bundle $\Sbundle':=G_{\overline0}\times_H S'$ to which elements of~$\Psi(\fg_{\overline 1})$ belong and which we could describe as the kernel of some fibrewise linear operator on~$\Sbundle$.
\end{Remark}

For the proof, we will make use of the technology for working with connections on homogeneous spaces presented in Appendix~\ref{sec:nomizu}, namely Nomizu maps and Wang's theorem (see Theorem~\ref{thm:Wang}). If we have a spin-lift $\tilde\varphi$ of the isotropy representation, then~${P=G_{\overline0}\times_H\Spin(V)}$ is a $G_{\overline0}$-invariant principal $\Spin(V)$-bundle and $\tilde\varphi$ is a map from the isotropy group to the structure group of $P$ -- to connect with the notation in Theorem~\ref{thm:Wang}, we have $G=G_{\overline 0}$, $K=\Spin(V)$ and $\psi=\tilde\varphi$. Then one can easily check that the Nomizu map~${L\colon \fg\to\fso(V)}$ of the Levi-Civita connection (considered as a $G_{\overline0}$-invariant principal connection on $F_{\rm SO}$) satisfies the conditions of Wang's theorem so also corresponds with a principal connection on $P$, which is of course the spin-lift of the Levi-Civita connection.

\begin{proof}
	By hypothesis, we have a metric Klein pair $(G_{\overline0},H,\eta)$. Then by Lemma~\ref{lemma:spin-lift-Ad}, there exists a spin lift of the isotropy representation of $(G_{\overline0},H,\eta)$ with corresponding homogeneous spin structure $P=G_{\overline0}\times_H\Spin(V)$ and spinor bundle $\Sbundle=G_{\overline0}\times_H S$ on $M=G_{\overline0}/H$. We identify $V$ with $T_oM\cong\fg/\fh$. Since $H$ is connected, the image of the spin-lift of the isotropy representation lies in $\Spin_0(V)$, whence $(M,g)$ is time-orientable and so we can construct a bundle Dirac current~${\kappa\colon \Odot^2\Sbundle\to TM}$ by Lemma~\ref{lemma:bundle-dirac-current-exist}.
	
Now let $[\mu_-]$ be the admissible Spencer cohomology class for $\fa=\Gr\fg\subseteq \fs$ corresponding to~$\fg$, so that there exists \smash{$\beta+\gamma\in\bigl(\cH^{2,2}\bigr)^\fh$} satisfying $i_*[\mu_-]=i^*[\beta+\gamma]\in\ssH^{2,2}(\fa_-;\fs)$. Recall that~${\beta\in\Hom(V\otimes S,S)\cong V^*\otimes\End S}$ is $\fh$-invariant and thus $H$-invariant, so by Frobenius reciprocity (see Theorem~\ref{thm:Frobenius-recipr}) it defines a unique $G_{\overline0}$-invariant 1-form with values in $\End\Sbundle$, although we will actually adopt a sign convention where we work with $\boldsymbol{\beta}\in\Omega^1(M;\End\Sbundle)^{G_{\overline0}}$ corresponding to~$-\beta$. We define the connection $D:=\nabla-\boldsymbol{\beta}$.

Since $G_{\overline 0}$ acts by isometries by construction and $\boldsymbol{\beta}$ is $G_{\overline 0}$-equivariant, for each $X\in\fg_{\overline 0}$ the fundamental vector field $\xi^M_X\in\fX(M)$ is Killing and preserves $\boldsymbol{\beta}$, so it lies in $\fV_D$. Recalling that the assignment $X\to\xi^M_X$ is a Lie algebra \emph{anti}-homomorphism $\fg_{\overline 0}\to \fX(M)$, we define $\Psi(X)=-\xi^M_X$. This gives us a Lie algebra morphism $\Psi\colon\fg_{\overline 0}\to\fV_D$.
	
	To construct the spinor fields associated to elements of $\fg_{\overline 1}=S'$ and show that they are $D$-parallel, we will exploit the homogeneous bundle structure of $S$. Let us define a left action of $H$ on $G_{\overline 0}$ by $h\cdot g=gh^{-1}$ where we use group multiplication in $G_{\overline 0}$ on the right-hand side. Then, using some standard associated bundle theory, we have identifications between spaces of sections and $H$-equivariant maps
	\[
		\fX(M) \cong C^\infty\qty(G_{\overline 0};V)^H,\qquad
		\fS \cong C^\infty\qty(G_{\overline 0};S)^H.
	\]
	For $X\in\fg_{\overline 0}$, the fundamental vector field $\xi^M_X$ corresponds to the map $G_{\overline 0}\to V\cong\fg_{\overline 0}/\fh$ given by $X\mapsto \overline{\Ad_{g^{-1}}X}$ (where the overline denotes the image under the quotient by $\fh$). Similarly, for $s\in \fg_{\overline 1}$ we define $\Psi(s)\in\fS$ as the spinor field associated to the map $G_{\overline 0}\to S$ defined by~${g\mapsto \Ad_{g^{-1}}s}$. We now use Wang's theorem and the Levi-Civita Nomizu map to show that these spinor fields are $D$-parallel. First, let us take a presentation of the form \eqref{eq:admiss-cocycle-brackets} for $\fg$; in particular, we have a (non-reductive) split $\fg_{\overline0}=\fh\oplus V$ and a map $\lambda\colon V\to\fso(V)$. Then, writing elements $X\in\fg_{\overline 0}$ in the form $X=A+v$ for $A\in\fh$, $v\in V$, we can follow the construction of the Nomizu map $L\colon \fg_{\overline 0}\to\fso(V)$ corresponding to the Levi-Civita connection (see Appendix~\ref{sec:nomizu}) to compute
$L(A+v) = A + \lambda(v)$.
	In particular, for $s\in S'$ we have $L(A+v)\cdot s = \comm{A+v}{s} - \beta(v,s)$, where we use the definition of the bracket on $\fg$. Since this is also the map corresponding to the spin-lift of the Levi-Civita connection $\nabla$ under Wang's theorem, we can compute the $H$-equivariant map $G_{\overline 0}\to S$ corresponding to the spinor field $\nabla_{\xi^M_X}\Psi(s)$ as
	\[
		g \mapsto (L(\Ad_{g^{-1}}X) - \ad_{\Ad_{g^{-1}}X})\qty(\Ad_{g^{-1}}s) = - \beta\qty(\overline{\Ad_{g^{-1}}X},\Ad_{g^{-1}}s),
	\]
	where the equality is due to the expression for $L$ above. But this is nothing but the map corresponding to $\boldsymbol{\beta}\bigl(\xi^M_X,\Psi(s)\bigr)$ (recalling our sign convention for $\boldsymbol{\beta}$). By homogeneity, the values of the fundamental vector fields span every tangent space, so we have shown that $\nabla\Psi(s)=\boldsymbol{\beta}\Psi(s)$, whence $\Psi(s)\in\fS_D$ as required. We now have a $\ZZ_2$-graded map $\Psi\colon\fg\to\fV_D\oplus\fS_D$ which can show satisfies the equations~\eqref{eq:reconstruct-brackets} by examining the Killing transport data at the basepoint (see Section~\ref{sec:alg-structure-killing}, in particular the explicit presentation of the brackets of the Killing superalgebra as a filtered deformation \eqref{eq:KSA-deformed-brackets}). This will also show that $\Psi$ is injective since no non-zero element of $\fg$ has trivial transport data.
	
	Finally, we note that if $D$ is admissible, $\Psi\colon\fg\hookrightarrow\fK_D$ is a Lie superalgebra embedding. The Spencer cocycle conditions for $\beta+\gamma$ ensure $\boldsymbol{\beta}$ satisfies conditions~(1) and~(2) of Definition~\ref{def:killing-spinor}. Condition~(3) ($\eL_{\kappa_\epsilon}\boldsymbol{\beta}=0$ for all $\epsilon\in\fS_D$) is not automatically satisfied, but observe that we have~$\eL_{\xi^M_X}\boldsymbol{\beta}=0$ for all $X\in\fg_{\overline 0}$ by $G_{\overline 0}$-invariance of $\boldsymbol{\beta}$; in particular, $\eL_{\kappa_\epsilon}\boldsymbol{\beta}=0$ holds for~${\epsilon\in\Psi(\fg_{\overline 1})}$.
\end{proof}

We note that we have glossed over some details of the calculations above, in particular when constructing the spinorial part of the map $\Psi$; these details are routine exercises in principal connections but somewhat cumbersome to write out in full. A more elegant way of dealing with them is through the superspace formalism introduced in \cite{Santi2010} which is used, somewhat implicitly, in the proof of \cite[Theorem~13]{Figueroa-OFarrill2017_1}, but which is beyond the scope of this work.
